% B3+ / ICML submission draft.
% Every number in this manuscript is produced from saved per-episode records by
% scripts/promising_directions/b3_plus_analysis.py and the table generators in
% paper/b3_plus/tables/. No value is hand-entered. Check the official ICML style
% requirements for the target year before formatting the final submission.
\documentclass{article}
\usepackage{booktabs,amsmath,amssymb,graphicx,hyperref,microtype}
\title{What Does It Cost to Keep a Certificate?\\
Label-Efficient Sequential Certification for Model Switching\\
under Budgeted and Delayed Feedback}
\author{}
\date{}
\begin{document}
\maketitle

\begin{abstract}
Deciding whether to replace a deployed model with a challenger requires paying for labels,
waiting for feedback, and absorbing a switching cost. Anytime-valid procedures answer this with a
sequential certificate, but they pay for that certificate in labels, and under a hard label budget
that price is itself a first-order cost. We show that a standard disagreement-gated paired
e-process keeps buying labels at its budget-paced rate long after the deployment decision is
settled: on one stream it promotes the challenger at event~244 and then purchases labels for the
remaining 1756 events. We introduce \emph{settled-decision acquisition}, which withholds queries
in the two states where no label can change the deployed action, and leaves pacing otherwise
untouched so it can only ever spend less. On four of five real-world case studies the resulting
policy reproduces the incumbent procedure's deployment decisions \emph{event for event} while
using 20--67\% fewer labels; on the fifth it alters behaviour and still improves on average. The
effect is invariant to the switching cost and exactly proportional to the label price, as a pure
label-saving mechanism must be, and it fails that scaling law on precisely the one dataset where
deployment behaviour changes. We do not claim better switching decisions: on one of our five cases
a fixed-sample binomial test remains the cheaper policy overall, and we report it alongside.
\end{abstract}

\section{Introduction}
A challenger model may be better than the incumbent in production, but acting on that possibility
is expensive. Labels must be purchased, feedback arrives late, and every switch of the deployed
model costs something. The natural tool is a sequential test that is valid at every stopping time,
so that promotion can happen as soon as the evidence warrants it.

Our starting point is an observation about what such a procedure spends. Optional stopping is
bought with labels, and the procedure keeps buying them for the whole episode. Yet once the
challenger has been promoted and continues to look better, the deployed action is not in question;
further labels buy only rollback insurance. Measured on real streams, this is where most of the
budget goes.

\paragraph{Contributions.}
\begin{itemize}
  \item We identify and quantify post-decision certificate spend in a disagreement-gated paired
  e-process, using saved per-episode records (Section~\ref{sec:method}).
  \item We propose settled-decision acquisition, which gates querying on whether a label could
  still change the deployed action. The evidence process, its propensity-matched bet cap, the
  promotion threshold and the rollback test are unchanged, and the rule can only withhold spend.
  \item We evaluate on five real-world case studies, report a falsifiable scaling law that the
  mechanism must satisfy, and show it holds on four cases and breaks on exactly the fifth, where
  its premise is independently violated (Sections~\ref{sec:results}--\ref{sec:ablation}).
\end{itemize}

\section{Related Work}
Contextual active model selection (CAMS) queries on a disagreement score and optimises regret with
the label budget as a hard constraint; it models no switching cost, does not use the remaining
horizon, does not address delay, and never commits to a deployed model, so it has no notion of
promotion or its price. Active, anytime-valid risk-controlling prediction sets give the closest
treatment of budgeted e-process-driven labelling: the policy maximises a lower bound on e-process
log-growth, driven by conditional risk variability rather than proximity to a decision boundary,
assumes labels are available immediately, and does not consider switching between candidate
models. Our acquisition rule is in one respect the opposite of that objective: it deliberately
declines to buy evidence once the evidence cannot change the action. We also compare against a
cost-aware promotion rule already present in our own codebase, which gates promotion on a lower
confidence bound times a declared persistence horizon; treating it as prior work rather than as a
baseline would overstate our novelty.

\section{Problem Formulation}
A chronological stream of $T$ events. A frozen incumbent $f_0$ and challenger $f_1$ produce $0/1$
losses $\ell_0[t],\ell_1[t]$. The deployed action $a_t\in\{0,1\}$ is chosen before $Y_t$ is
observed. A label costs $c_q$, at most $B$ may be requested per episode, and a requested label is
released $d$ events later. Each switch costs $c_s$. The realised net cost is
\[
  C_\pi=\sum_{t=1}^{T}\ell_{a_t}[t]+c_qN_q+c_sN_s .
\]
Each event proceeds in three stages. \textbf{Issue}: choose a query probability $q_t$ and freeze a
betting fraction $\lambda_t$, both measurable with respect to information released strictly before
event $t-d$; only disagreement events are queried, and $q_t\ge p_{\min}$ while budget remains.
\textbf{Release}: the label of $s=t-d$ becomes usable and its importance-weighted paired increment
$x_s=(\ell_0[s]-\ell_1[s])/q_s$ enters the e-process with the $\lambda_s$ frozen at its own issue
time. \textbf{Decide}: promote when $\log W_t\ge\log(1/\alpha)$, roll back by the mirrored test,
effective from $t+1$.

\section{Method}\label{sec:method}
\paragraph{Where the budget goes.} With the incumbent procedure, promotion occurs at event $244$
on \texttt{spam}, $425$ on \texttt{rialto} and $988$ on \texttt{weather}, each at promotion rate
$1.00$, after which querying continues at the paced rate for the remainder of the episode.

\paragraph{Settled-decision acquisition.} Two states admit no label that could change the deployed
action: \textsc{hopeless}, not yet promoted with the challenger certified worse; and
\textsc{settled}, already promoted with the recent discounted paired advantage still favouring the
challenger. Writing $r_t$ for a multiplier on the paced rate,
\[
  q_t=\mathrm{clip}\!\left(r_t\cdot\frac{\mathrm{rem}_t}{\widehat{\text{rate}}_t\,(T-t)},\;
  p_{\min},\;1\right),\qquad
  r_t=\begin{cases}
    0.25 & \textsc{hopeless}\\
    0 & \textsc{settled}\\
    1 & \text{otherwise.}
  \end{cases}
\]
Since $r_t\le 1$ the rule can only withhold spend. A settled multiplier of $0$ still queries at
$p_{\min}$ because $q_t$ is clipped there, so positive propensity and rollback monitoring both
survive. Post-promotion the margin grids are frozen, so the settled test uses the recent
discounted advantage, which remains defined.

\paragraph{What is and is not guaranteed.} The margin-null e-processes have non-negative wealth
factors under the propensity-matched cap, issued propensities are measurable with respect to
released information and bounded below by $p_{\min}$, and the budget is pathwise feasible, so the
certificate is unaffected: acquisition steers spending, not validity. We claim no guarantee about
rollback latency; reduced post-promotion monitoring detects a reversal later than the incumbent
would, and we characterise that empirically rather than theoretically.

\section{Experimental Setup}
Five real-world streams: \texttt{spam}, \texttt{rialto}, \texttt{weather}, \texttt{openml\_ldpa}
and \texttt{uci\_occupancy}. Primary configuration $T=2000$, $B=200$, $c_q=0.05$, $c_s=5$, with
dataset-specific imposed delays declared before any locked window was read, and $100$
query-randomisation seeds shared across policies so every comparison is paired on episodes and
seeds. Comparators: $P_0$ never switch; P3b the incumbent certified procedure; P6 a fixed-sample
exact binomial test on released disagreement labels; P2 incumbent-alarm switching. $P_1$, always
deploying the challenger, is a hindsight reference and never a deployable baseline. Model pairs
are selected on train/dev windows only.

\paragraph{Selection disclosure.} The five cases are a \emph{selected} set. Thirty-six datasets
were considered, twenty-nine reached a screening verdict, eleven reached locked-window evaluation,
and five were selected. The datasets named in our pre-registered queue produced no successes; both
new cases come from candidates added after that queue was exhausted. These figures are not an
unbiased success rate, and the method was not evaluated only on these five.

\section{Main Results}\label{sec:results}
\begin{table}[t]\centering\small
\begin{tabular}{llrrrrrrr}
\toprule
Dataset & Domain & $K$ & $P_0$ & P3b & P6 & \textbf{P7s} & P7s$-$P3b (95\% CI) & labels \\
\midrule
spam$^{\dagger}$ & Email spam & 1 & 849.00 & 254.09 & 315.95 & \textbf{247.43} & $-6.66$ $[-6.77,-6.55]$ & 200$\to$67 \\
rialto & Image stream & 16 & 202.06 & 103.73 & 122.91 & \textbf{98.42} & $-5.31$ $[-6.00,-4.58]$ & 159$\to$53 \\
weather$^{\dagger}$ & Environmental & 3 & 504.00 & 475.77 & 465.56 & \textbf{472.57} & $-3.20$ $[-3.32,-3.08]$ & 178$\to$114 \\
ldpa & Wearable HAR & 32 & 663.00 & 483.25 & 506.51 & \textbf{478.82} & $-4.43$ $[-7.05,-2.58]$ & 181$\to$115 \\
occupancy$^{\dagger}$ & Occupancy & 4 & 481.25 & 177.63 & 236.45 & \textbf{176.62} & $-1.01$ $[-1.03,-0.99]$ & 101$\to$81 \\
\bottomrule
\end{tabular}

\caption{Net cost per episode at the primary configuration. $\dagger$ marks cases whose interval
is a seed-level bootstrap ($K<5$), which measures query-randomisation variance and not episode
variance. Only \texttt{rialto} and \texttt{openml\_ldpa} are properly powered.}
\end{table}

The method improves on the incumbent on all five cases with every confidence interval below zero,
and reduces labels by 20--67\%. The improvement is not uniform across comparators: on
\texttt{weather} the fixed-sample test P6 remains cheaper than both the incumbent and our method,
and P6 also beats the incumbent on the four streams that failed our selection condition. We
therefore claim an improvement \emph{on the certified procedure}, in the regime where that family
is appropriate, and not that our policy is the best available.

\section{Mechanism and Ablations}\label{sec:ablation}
Decomposing the paired difference into deployment loss, query cost and switch cost leaves the
entire effect in the query term on four of five cases, where deployment loss is bitwise identical
per episode, promotion times coincide and rollbacks are zero for both policies. The method
therefore does not improve decisions; it removes waste.

\paragraph{A falsifiable scaling law.} If the whole gain is $c_qN_q$, the paired difference must be
invariant to $c_s$ and exactly proportional to $c_q$. Both predictions hold to displayed precision
on \texttt{spam}, \texttt{rialto}, \texttt{weather} and \texttt{uci\_occupancy}, and both fail on
\texttt{openml\_ldpa}, the one case where deployment behaviour changes. A fitted effect would not
reproduce an exact scaling law on four datasets and break on precisely the dataset independently
shown to violate its premise.

\begin{table}[t]\centering\small
\begin{tabular}{lrrrrr}
\toprule
Dataset & S0 & S1 ($B{=}400$) & S2 ($c_s{=}20$) & S3 ($2\times$ delay) & S4 ($c_q{=}0.10$) \\
\midrule
spam & $-6.66$ & $-15.62$ & $-6.66$ & $-6.18$ & $-13.32$ \\
rialto & $-5.31$ & $-5.43$ & $-5.31$ & $-4.14$ & $-10.62$ \\
weather & $-3.20$ & $-5.88$ & $-3.20$ & $-3.17$ & $-6.40$ \\
openml\_ldpa & $-4.43$ & $-11.48$ & $-5.29$ & \textemdash & $-7.73$ \\
uci\_occupancy & $-1.01$ & $-2.79$ & $-1.01$ & \textemdash & $-2.02$ \\
\bottomrule
\end{tabular}

\caption{Paired difference by configuration. Em-dashes mark configurations not run rather than
interpolated values.}
\end{table}

The benefit scales with \emph{unspent budget the incumbent would otherwise consume}, not with the
budget granted: raising $B$ to $400$ doubles to triples the gain on three cases, but leaves
\texttt{rialto} essentially unchanged, because its disagreement rate caps queries near $165$ so
the larger budget never binds.

\section{Discussion and Limitations}
Three limitations are load-bearing. First, power: three of five cases contribute fewer than five
episodes, so their intervals are seed-level bootstraps; a second disjoint seed batch reproduces
every effect to within 1--3\% but does not create episode-level power. Second, the reversal risk:
reduced post-promotion monitoring detects reversal later, and while \texttt{openml\_ldpa} exercises
reversal on real data and our method rolls back \emph{less} often there with better mean deployment
loss, it costs more in 2 of 32 episodes, and we have no real stream where reduced monitoring is
worse on average. Third, \texttt{openml\_ldpa}'s event order is inferred from undocumented
columns rather than documented, and its non-stationarity is largely a change of recorded subject
between blocks rather than drift within one deployment.

\section{Conclusion}
Keeping a sequential certificate valid costs labels, and much of that cost is incurred after the
deployment decision has already been made. Withholding queries in the states where no label can
change the deployed action recovers that spend, on four of five real case studies with deployment
behaviour identical event for event. The result is an efficiency claim, not a better decision
rule, and we have reported where a simpler procedure still wins.

\end{document}
